\section{Discussion and Future Work}\label{sec:discussion}

The calibration and scenario results validated the core claim of this study: a modified Heston process with a state-dependent $\theta$-function produced realistic IV surfaces from structural return dynamics alone, and the same decomposition of $\theta$ into a per-ticker level and a shape function $\psi$ accommodated a hierarchy of representations spanning a parsimonious analytic form, a globally shared neural surrogate, and a sector-specific neural surrogate. The single-DTE parametric smile captured the essential features of the NVDA IV smile, and the smile shape generalized across semiconductor names with only a scalar level adjustment; the multi-expiration calibration extended this to the full NVDA IV surface with the five-parameter $\psi$ function and the quadratic DTE term $\beta_5$ that captured the U-shaped ATM term structure. The equilibrium initialization $v_0 = \theta(t\!=\!0)$ eliminated $v_0$ as a free parameter and produced the smile, skew, and term structure by construction; this was not merely a computational convenience but a structural constraint that ensured the initial IV surface was fully determined by the $\theta$-function, consistent with the view that IV surfaces reflect the conditional distribution of returns under the risk-neutral measure~\cite{cont2002}. The $\psi$-function played an analogous role to the SVI parameterization~\cite{gatheral2004} in that it provided a parsimonious representation of the smile and term structure, but with a critical difference: the parameters entered a generative model rather than a static interpolation, so the calibrated surface could evolve dynamically as the underlying HMM state changed.

Scaling the calibration to the 31-ticker two-day ladder revealed the limits of the parametric form as a universal shape function and motivated the hierarchy of neural representations. A single global $\psi$ fit across the six sectors achieved an overall RMSE of 9.36\% IV, dominated by systematic misfits in sectors whose smile geometry departed from the semiconductor baseline (healthcare at 10.78\%, technology at 11.65\%). Replacing the parametric form with a globally shared neural $\psi$ reduced the overall RMSE to 8.03\% without changing the sharing granularity, confirming that part of the gap was a functional-form bottleneck in the log-polynomial representation. Loosening the sharing to one network per sector reduced the overall RMSE further to 6.76\%, with the sector NN strictly dominating the shared NN in every sector, which established that both axes (capacity and granularity) contributed independently. The remaining error was concentrated in two structural sources: within-sector heterogeneity in technology and healthcare, where specific tickers (INTC, AVGO, MU, PFE, ABBV) carried idiosyncratic smile shapes not expressible by a shared sector-level $\psi$, and small-sample effects for tickers with fewer than 200 observations (PFE at 101 observations). Both sources are addressable through the same decomposition: the per-ticker roadmap relaxes shape sharing to the ticker level for the data-rich names, and additional capture days directly reduce small-sample variance for the thin names.

Extending the corpus to eight capture days and running a temporal holdout (Section~\ref{sec:earnings_holdout}) reframed the generalization question more sharply than the in-sample RMSE alone could. The pooled-control sector NN trained on six days and evaluated on the next two carried a $4.99\%$ IV generalization gap ($8.05\%$ train, $13.03\%$ test), which on its face suggested a substantial overfit. Excluding observations within $\pm 3$ days of any same-sector earnings print from both train and test eliminated that gap entirely ($8.02\%$ train, $7.96\%$ test, $-0.06\%$ generalization gap), revealing that the architecture itself generalized faithfully out of sample and that the entire $4.99\%$ gap was attributable to scheduled corporate events embedded in the holdout window. Adding earnings-distance and same-sector peer-distance features to the shape network closed roughly one fifth of the gap on the original train and test rows ($11.98\%$ test, $+4.36\%$ generalization gap). The per-ticker decomposition exposed the mechanism: the peer-coupling feature carried the IV-spillover signal, so QCOM (own print six days after the window) and AMD (twelve days after) saw $\sim 4\%$ IV reductions in test error from the INTC print on a holdout day, and the same mechanism cut META, MSFT, and GOOG test RMSE by 2 to $4\%$ IV. INTC itself was the only Tech name whose error rose under the earnings-aware fit, driven by an outsized $+2192\%$ EPS surprise that the calendar-based indicator could not predict. The boundary was therefore sharp: the indicator captured the anticipatory components of event-day IV behavior (pre-print expansion on the printer, sympathy on same-sector peers), which were a function of the calendar alone, but it carried no information about the post-event realization when consensus was wrong. The non-earnings baseline of $7.96\%$ test RMSE is the right reference for what the current architecture can do under faithful temporal generalization; the residual gap to the pooled $13.03\%$ quantifies the part of the surface that requires a richer event mechanism than the calendar.

Several limitations constrained the scope of the current conclusions and pointed directly to the natural extensions that followed from them. The cross-sector partition used conventional GICS-style groupings; whether alternative groupings (mega-cap versus semiconductor within technology; pharma versus biotech within healthcare) improve the sector-NN fit without fragmenting the corpus is an empirical question that additional data will clarify. The forward scenario analysis (Section~\ref{sec:scenarios}) confirmed that the calibrated parameters produced IV paths centered on observed market levels with realistic stochastic fluctuations, but the validation was limited to a single ticker and a single market regime; broader validation across sectors and volatility environments, and a re-run of the scenarios with the sector-NN $\psi$ for each represented sector, would strengthen the generalizability claim. The earnings-aware fit recovered the anticipatory portion of the event-day signal but left the surprise-driven portion uncaptured. A natural extension restored the per-ticker tail term that the simulation pipeline already supports for forward state evolution and that the calibration omitted: replacing the operational $\theta_i$ with $\theta_i \cdot (1 + \gamma_i M_i^{(t)})$, where $M_i^{(t)}$ was the JumpHMM tail-state probability for ticker $i$ at time $t$ and $\gamma_i$ was a learned per-ticker tail sensitivity, would absorb idiosyncratic IV inflation from any source whose return-side imprint already drove the HMM into a tail state, including but not limited to earnings surprises. Restoring this term required state-conditioned histories rather than observation-day surfaces, so it sat naturally at the simulation interface: the calibrated $\psi_{\text{NN}}$ remains the IV shape function, the calendar-based earnings inputs $(e,\,e^{\text{peer}})$ remain the scheduled-event indicator, and $\gamma_i M_i$ becomes the unscheduled tail mechanism populated from the JumpHMM at simulation time. The natural reference point for evaluating either of these extensions is the non-event $7.96\%$ test RMSE established in Section~\ref{sec:earnings_holdout}, against which any further improvement would represent a genuine reduction in the residual event signal rather than a recovery of cross-sectional fit already accounted for by the sector $\psi$. Beyond these specific limitations, the primary intended use of this framework was generating realistic synthetic training data for machine-learning models in finance, with applications ranging from training option price predictors and learning delta-hedging policies via reinforcement learning~\cite{buehler2019} to augmenting scarce real option data for transfer learning. The framework naturally supported computation of Greeks ($\Delta$, $\Gamma$, $\mathcal{V}$, $\Theta$) via finite differences on the CRR tree, and including these in the synthetic dataset would support hedging applications. Finally, the current mood mechanism created correlated IV dynamics through a shared stress signal; a richer model could introduce ticker-pair IV correlations through the copula structure, producing a full cross-asset IV surface that respected observed correlation patterns between individual-name and index volatility.
